% =====================================================================
%  flux_dose.tex  --  ICRP-116 neutron fluence-to-effective-dose
%  conversion coefficients (the DE2/DF2 data of Table~\ref{tab:icrp}),
%  plotted on logarithmic axes.  \input from chapters/4-simulation.tex.
% =====================================================================
\begin{figure}[H]
  \centering
  \begin{tikzpicture}
    \begin{axis}[
      width=0.95\textwidth, height=7.35cm,
      xmode=log, ymode=log,
      xlabel={neutron energy, $E$~[MeV]},
      ylabel={effective dose, $h$~[pSv\,cm$^2$]},
      xmin=5e-10, xmax=1e2,
      ymin=1,     ymax=1000,
      axis lines=left,
      axis line style={line width=0.4pt, -{Triangle[length=2.6mm,width=2.4mm]}},
      xtick={1e-9,1e-8,1e-7,1e-6,1e-5,1e-4,1e-3,1e-2,1e-1,1e0,1e1,1e2},
      ytick={1,3,10,30,100,300,1000},
      yticklabels={1,3,10,30,100,300,1000},
      minor ytick={2,4,5,6,7,8,9,20,40,50,60,70,80,90,200,400,500,600,700,800,900},
      yminorticks=true,
      grid=major,
      grid style={gray!30},
      tick label style={font=\small},
      label style={font=\small},
    ]
      % vertical bin dividers from the E2 card (0.5e-6, 0.5, 20 MeV):
      % thermal | epithermal | fast.  Drawn first so the curve sits on top;
      % they span the full plot height (ymin=1 .. ymax=1000).
      \draw[densely dashed, black!70] (axis cs:0.5e-6,1) -- (axis cs:0.5e-6,1000);
      \draw[densely dashed, black!70] (axis cs:0.5,1)    -- (axis cs:0.5,1000);
      \draw[densely dashed, black!70] (axis cs:20,1)     -- (axis cs:20,1000);
      % curve split into the three energy bins, each in its own colour; the
      % bin-edge points (5e-7 and 5e-1) are shared so the line stays continuous.
      % thermal segment (orange): E <= 0.5e-6 MeV
      \addplot[orange, thick, mark=*, mark size=1.1pt, mark options={fill=orange}]
        coordinates {
          (1e-9,3.09) (1e-8,3.55) (2.5e-8,4.00) (1e-7,5.20) (2e-7,5.87) (5e-7,6.59)
        };
      % epithermal segment (blue): 0.5e-6 < E <= 0.5 MeV
      \addplot[blue, thick, mark=*, mark size=1.1pt, mark options={fill=blue}]
        coordinates {
          (5e-7,6.59)   (1e-6,7.03)   (2e-6,7.39)   (5e-6,7.71)   (1e-5,7.82)
          (2e-5,7.84)   (5e-5,7.82)   (1e-4,7.79)   (2e-4,7.73)   (5e-4,7.54)
          (1e-3,7.54)   (2e-3,7.61)   (5e-3,7.97)   (1e-2,9.11)   (2e-2,12.2)
          (3e-2,15.7)   (5e-2,23.0)   (7e-2,30.6)   (1e-1,41.9)   (1.5e-1,60.6)
          (2e-1,78.8)   (3e-1,114)    (5e-1,177)
        };
      % fast segment (green): 0.5 < E <= 20 MeV
      \addplot[green!50!black, thick, mark=*, mark size=1.1pt,
               mark options={fill=green!50!black}]
        coordinates {
          (5e-1,177)  (7e-1,232) (9e-1,279) (1,301)  (1.2,330) (1.5,365) (2,407)
          (3,458)     (4,483)    (5,494)    (6,498)  (7,499)   (8,499)   (9,500)
          (10,500)    (12,499)   (14,495)   (15,493) (16,490)  (18,484)  (20,477)
        };
      % section labels, centred horizontally in each bin and all at the same
      % height: the vertical centre of the plot (geom. mean of 1..1000 = 31.6).
      % white fill masks the nearby y=30 gridline behind the text.
      \node[font=\small, orange,         fill=white, inner sep=2pt] at (axis cs:1.5e-8,31.6) {thermal};
      \node[font=\small, blue,           fill=white, inner sep=2pt] at (axis cs:5e-4,31.6)   {epithermal};
      \node[font=\small, green!50!black, fill=white, inner sep=2pt] at (axis cs:3.2,31.6)    {fast};
      % close the frame: plain top and right borders (no ticks), matching the axes
      \draw[black, line width=0.4pt]
        (rel axis cs:0,1) -- (rel axis cs:1,1) -- (rel axis cs:1,0);
    \end{axis}
  \end{tikzpicture}
  \caption{The conversion coefficients of table~\ref{tab:icrp} plotted on logarithmic axes.}
  \label{fig:icrp}
\end{figure}
